\section{Integer and mixed harmonic specializations}
\label{dhr:sec:harmonic}

\subsection{A rational spectral partition at every particle number}
Take $a=c=1$, so $b_0=2r$, and choose
\begin{equation}\label{dhr:eq:integer-normalization}
 L(u)=\frac{\Gamma(2r+u)}{\Gamma(2r)}.
\end{equation}
At zero the particle weight is rational:
\[
 w_n(1,1;2r)=\frac1{(n+1)_{2r-1}}.
\]
Denote the normalized meromorphic partition by
$F_r(u)=L(u)^rZ_r(1,1;2r+u)$.

\begin{proposition}[Rationality of the normalized partition]
\label{dhr:prop:rational-partition}
For every $r\geq1$,
\begin{equation}\label{dhr:eq:rational-partition}
 \boxed{F_r(u)=
 \frac{\displaystyle\prod_{j=0}^{r-1}(j!)^3}{[(2r-1)!]^r}
 \prod_{k=0}^{r-1}\frac{(u+r+k)_{r-k}}{(u+k)_r}.}
\end{equation}
Its residue $A_0$ at zero is a positive rational number. The hard-cutoff
constant is
\begin{equation}\label{dhr:eq:integer-cutoff}
 A_0\{\gamma-3H_{r-1}\}.
\end{equation}
\end{proposition}
\begin{proof}
For $a=c=1$, the $j$th product factor in
\eqref{dhr:eq:gauss-product} is
\[
 (j!)^3\frac{\Gamma(r-1+u-j)}{\Gamma(2r-1+u-j)^2}.
\]
Pair one factor $\Gamma(2r+u)$ from $L^r$ with each of these,
use Gamma recurrences, and replace $j$ by $r-1-k$.
This gives \eqref{dhr:eq:rational-partition}. Finally,
$\chi=2\psi(1)-3\psi(r)=\gamma-3H_{r-1}$, and
Corollary~\ref{dhr:cor:cutoff} applies.
\end{proof}

The first partitions are
\begin{align}
 F_1(u)&=\frac{u+1}{u},\notag\\
 F_2(u)&=\frac{(u+3)^2}{36u(u+1)^2},\notag\\
 F_3(u)&=\frac{(u+4)(u+5)^3}
 {216000u(u+1)^2(u+2)^3(u+3)},\notag\\
 F_4(u)&=\frac{(u+6)^2(u+7)^4}
 {373403520000u(u+1)^2(u+2)^3(u+3)^4(u+4)^2}.
 \label{dhr:eq:small-partitions}
\end{align}
The corresponding residues are
\begin{equation}\label{dhr:eq:small-residues}
 A_0=1,\quad\frac14,\quad\frac1{10368},\quad
 \frac1{44789760000}\qquad(r=1,2,3,4).
\end{equation}
Although $F_r$ is rational, the individual normal-convergence amplitude
$A(u)$ is generally not rational. Its polygamma derivatives are exactly
what permit Stieltjes constants to appear in ordinary subtracted sums.

For calculation, the logarithmic derivatives are
\begin{align}
 (\log A)^{(k)}(0)
 &=r\psi^{(k-1)}(2r)
 +\sum_{j=0}^{r-2}\left\{\psi^{(k-1)}(r-j)
                    -2\psi^{(k-1)}(2r-1-j)\right\},
 \label{dhr:eq:integer-A-jets}\\
 \ell_k&=(-1)^k(k-1)!
       \{3H_{r-1}^{(k)}-\zeta(k)\},\qquad k\geq2.
 \label{dhr:eq:integer-E-jets}
\end{align}
Together with $\ell_1=\gamma-3H_{r-1}$, these are finite exact
coefficient recipes for every row of Theorem~\ref{dhr:thm:jets}.

\subsection{The derivatives as actual harmonic polynomials}
For a configuration $\bm n=(n_1,\ldots,n_r)$ define
\begin{equation}\label{dhr:eq:harmonic-log-jets}
 v_k(\bm n)=(-1)^k(k-1)!
  \sum_{i=1}^r\left\{H_{n_i+2r-1}^{(k)}-H_{2r-1}^{(k)}\right\}.
\end{equation}
Logarithmic differentiation of the normalized weight gives
\begin{equation}\label{dhr:eq:K-harmonic}
 K_n^{(k)}(0)=
 \sum_{0\leq n_1<\cdots<n_{r-1}<n_r=n}
 \frac{\Delta(\bm n)^2\,
  \Bell_k(v_1(\bm n),\ldots,v_k(\bm n))}
 {\displaystyle\prod_{i=1}^r(n_i+1)_{2r-1}}.
\end{equation}
This formula has no spectral derivatives left unevaluated. To see it,
for a single particle differentiate
$\log\Gamma(2r+u)-\log\Gamma(n+2r+u)$.
The polygamma shift formula gives the corresponding summand in
\eqref{dhr:eq:harmonic-log-jets}; the Bell identity
\eqref{dhr:eq:bell} then differentiates the product. Substitution in
\eqref{dhr:eq:all-jets} proves an explicit all-order family of
harmonic--logarithm--Stieltjes identities.

For a concrete first-order member, put
\begin{equation}\label{dhr:eq:W2}
 W_{m,n}=\frac{(n-m)^2}{(m+1)_3(n+1)_3},\qquad 0\leq m<n.
\end{equation}
At $r=2$, $A_0=1/4$, $A_1=(5/3-\gamma)/4$, and
$e_2=(\gamma-3)^2+3-\zeta(2)$. Thus
\begin{align}
 \sum_{n\geq0}\bigg[&\sum_{m=0}^{n-1}W_{m,n}
       (2H_3-H_{m+3}-H_{n+3})\notag\\
 &-\frac{5/3-\gamma-\log(n+1)}{4(n+1)}\bigg]
 =\frac14+\frac{\gamma^2-\zeta(2)}8+\frac{\gamma_1}{4}.
 \label{dhr:eq:explicit-stieltjes}
\end{align}
Here $\gamma_1=\gamma_1(1)$, and the empty inner sum at $n=0$ is zero.
The zero-order companion is
\begin{equation}\label{dhr:eq:explicit-stieltjes-zero}
 \sum_{n\geq0}\left\{\sum_{m<n}W_{m,n}-\frac1{4(n+1)}\right\}
 =-\frac34.
\end{equation}
These are ordinary convergent bracketed series, not a declaration of
finite-part integration in disguise.

The same example makes the normalization distinction numerical and exact:
\begin{equation}\label{dhr:eq:normalization-example}
 \FP_{u=0}F_2(u)=-\frac13,\qquad
 \lim_{N\to\infty}\left\{\sum_{m<n\leq N}W_{m,n}
                              -\frac14\log N\right\}
 =\frac{\gamma-3}{4}.
\end{equation}
Their difference is $A_1=(5/3-\gamma)/4$, as required. With $L=1$
the spectral finite part instead equals $(2\gamma-5)/4$; the hard
cutoff is unchanged.

\subsection{Ordinary all-order mixed harmonic sums}
A convergent, rather than resonant, specialization produces another
useful family. For an integer $B>2r$, define
\begin{equation}\label{dhr:eq:mixed-generator}
 \mathcal H_{r,B}(p,q)=
 \frac{Z_r(1+p,1+q;B)}{[\Gamma(1+p)\Gamma(1+q)]^r}.
\end{equation}
Near $p=q=0$, it is both a normally convergent sum and the explicit
Gamma product
\begin{align}
 \mathcal H_{r,B}(p,q)
 =\prod_{j=0}^{r-1}
 \frac{j!\,\Gamma(1+p+j)\Gamma(1+q+j)
              \Gamma(B-r-1-j-p-q)}
 {\Gamma(1+p)\Gamma(1+q)
              \Gamma(B-1-j-p)\Gamma(B-1-j-q)}.
 \label{dhr:eq:mixed-product}
\end{align}
The convergence region includes a neighborhood with
$\Rea(p+q)<B-2r$.

\begin{theorem}[Many-variable mixed harmonic identities]
\label{dhr:thm:mixed}
For a configuration $\bm n$, let
\[
 U_j(\bm n)=(-1)^{j-1}(j-1)!\sum_{i=1}^rH_{n_i}^{(j)},\qquad
 W_B(\bm n)=\frac{\Delta(\bm n)^2}{\prod_i(n_i+1)_{B-1}}.
\]
For all $k,l\geq0$,
\begin{align}
 &\sum_{n_1<\cdots<n_r}W_B(\bm n)
 \Bell_k(U_1,\ldots,U_k)\Bell_l(U_1,\ldots,U_l)\notag\\
 &\hspace{35mm}=\partial_p^k\partial_q^l
               \mathcal H_{r,B}(0,0).
 \label{dhr:eq:mixed-harmonic}
\end{align}
Every pure row ($k=0$ or $l=0$) is rational. Every mixed row lies
in $\Q[\zeta(2),\ldots,\zeta(k+l)]$; for zero or first total order
this means simply $\Q$.
\end{theorem}
\begin{proof}
Relative to its value at $p=q=0$, the summand is multiplied by
\[
 \prod_{i=1}^r\prod_{j=1}^{n_i}(1+p/j)(1+q/j).
\]
Its logarithmic derivatives in either variable are the $U_j$.
The Bell formula and local normal convergence prove
\eqref{dhr:eq:mixed-harmonic}. The pure generating function is
itself rational:
\begin{equation}\label{dhr:eq:pure-rational}
 \mathcal H_{r,B}(p,0)=
 \prod_{j=0}^{r-1}
 \frac{(j!)^2(1+p)_j}
 {\Gamma(B-1-j)(B-r-1-j-p)_r}.
\end{equation}
This follows by Gamma recurrences in \eqref{dhr:eq:mixed-product}.
It proves rationality of every pure derivative, not just the displayed
low orders. For derivatives involving both $p$ and $q$, only
$\log\Gamma(B-r-1-j-p-q)$ contributes to a mixed derivative
of the logarithm. Its derivatives of total order $h\geq2$ are
rational combinations of $\zeta(h)$ and finite generalized harmonic
numbers. Pure logarithmic derivatives are rational by
\eqref{dhr:eq:pure-rational}. Exponentiating with finite multivariate
Bell coefficients proves the asserted zeta algebra.
\end{proof}

For $r=2,B=6$, the product simplifies to
\begin{equation}\label{dhr:eq:H26}
 \mathcal H_{2,6}(p,q)=
 \frac{(1+p)(1+q)\Gamma(3-p-q)\Gamma(2-p-q)}
 {\Gamma(5-p)\Gamma(5-q)\Gamma(4-p)\Gamma(4-q)}.
\end{equation}
Set $V_{m,n}=(n-m)^2/[(m+1)_5(n+1)_5]$ and sum below over
$0\leq m<n$. Direct differentiation yields
\begin{align}
 \sum V_{m,n}&=\frac1{10368},\notag\\
 \sum V_{m,n}(H_m+H_n)&=\frac{29}{12\cdot10368},\notag\\
 \sum V_{m,n}\{(H_m+H_n)^2-H_m^{(2)}-H_n^{(2)}\}
 &=\frac{43}{8\cdot10368},\notag\\
 \sum V_{m,n}(H_m+H_n)^2
 &=\frac{517/144+2\zeta(2)}{10368},\notag\\
 \sum V_{m,n}\{H_m^{(2)}+H_n^{(2)}\}
 &=\frac{2\zeta(2)-257/144}{10368}.
 \label{dhr:eq:ordinary-examples}
\end{align}
For example, at zero the first logarithmic $p$-derivative is $29/12$,
the second pure logarithmic derivative is $-67/144$, and the mixed
logarithmic derivative is $2\zeta(2)-9/4$. These three numbers prove
all four nonconstant rows. They illustrate why a pure Bell combination
can be rational while the square of the ordinary harmonic sum is not.
